% TOP_PROBLEM: {"id": 3147, "title": "Ramsey properties of Cayley graphs", "category_id": 3, "category": {"id": 3, "name": "graph_theory", "display_name": "Graph Theory", "description": "Problems involving graphs, networks, and their properties.", "slug": "graph-theory", "order_index": 3, "created_at": "2026-07-31T15:26:25.671Z"}, "status": "open", "problem_number": "OPG-372", "set_id": 12, "statusQualification": "Finite groups of order at least two and identity-free connection sets make the asymptotic simple-graph formulation explicit."}
% FIRST_POSTED: 2026-10-01
% ATTEMPT_STATUS: unresolved
% UnsolvedMath ID 3147; source catalogue code OPG-372
% !TeX program = lualatex
% Research attempt by GPT-6 Astra Ultra; no claim of a new complete solution.
\documentclass[11pt,a4paper]{article}
\usepackage[margin=25mm]{geometry}
\usepackage{fontspec}
\setmainfont{lmroman10-regular.otf}[BoldFont=lmroman10-bold.otf,
 ItalicFont=lmroman10-italic.otf,BoldItalicFont=lmroman10-bolditalic.otf]
\usepackage{amsmath,amssymb,mathtools}
\usepackage{unicode-math}
\setmathfont{latinmodern-math.otf}
\AtBeginDocument{\let\setminus\smallsetminus}
\usepackage{xurl,hyperref}
\hypersetup{colorlinks=true,urlcolor=blue}
\setcounter{secnumdepth}{0}
\setlength{\parindent}{0pt}
\setlength{\parskip}{5pt}
\setlength{\emergencystretch}{3em}
\begin{document}
\section*{Ramsey properties of Cayley graphs}\label{3147}
{\small UnsolvedMath ID 3147 · OPG-372 · Graph Theory · OpenGarden}
\subsection{Definitions and mathematical statement}
Let $A$ be a finite abelian group, written additively, with
$N=|A|\geq2$ and identity $0$. For a set
$S\subseteq A\setminus\{0\}$ satisfying $S=-S$, define the simple
undirected Cayley graph $\operatorname{Cay}(A,S)$ on vertex set $A$
by joining distinct $x,y$ exactly when $y-x\in S$. Symmetry of $S$
makes adjacency independent of the order of the two vertices.
The set $S$ is not required to generate $A$.

A clique is a vertex set all of whose distinct pairs are adjacent;
an independent set has no adjacent pair. Their maximum sizes are
denoted $\omega(G)$ and $\alpha(G)$, respectively. The conjecture
asks whether there is an absolute real constant $c>0$ such that
\[
 \forall A\quad\exists S=-S\subseteq A\setminus\{0\}:
 \max\{\omega(\operatorname{Cay}(A,S)),
        \alpha(\operatorname{Cay}(A,S))\}\leq c\log N.
\]
Here $\log$ is the natural logarithm; changing its base only changes
the permitted constant. The same $c$ must work for groups of every
order and every isomorphism type, whereas $S$ may depend on the
group. No probabilistic construction is mandated. The restriction
$N\geq2$ removes the zero-logarithm boundary case and leaves the
asymptotic problem unchanged. Both homogeneous-set bounds must hold
for the same chosen Cayley graph.
\subsection{Short English statement}
Can every finite abelian group support a Cayley graph with neither a clique nor an independent set larger than a constant times the logarithm of its order?
\subsection{Research attempt}
\paragraph{Scope and process.}
This is a GPT-6 Astra Ultra research attempt dated 1 October 2026, using
primary-source browsing and elementary mathematical checks. The canonical
definition above is retained. Results below are restricted results or
reductions, not a claim of a new solution to the entire catalogue question.
No formal proof assistant or external mathematical referee checked this text.


\paragraph{Literature, conventions and plan.}
The question in [S2] is an existence problem for difference Cayley graphs.
Green's primary paper [R1] concerns random Cayley \emph{sum} graphs;
on groups of exponent two, sums and differences coincide. It shows
that the simplest random model there need not achieve the desired
logarithmic clique bound. Thus a random-edge heuristic cannot settle
the question for every abelian group. We instead prove a weaker uniform
existence bound by controlling all character eigenvalues, and locate
exactly where that bound loses the required scale.

\paragraph{An elementary random spectral construction.}
Partition $A\setminus\{0\}$ into inverse orbits $O=\{s,-s\}$, allowing
singleton orbits. Include each entire orbit independently in $S$ with
probability $1/2$. The graph is automatically simple and undirected.
For a character $\chi:A\to\mathbb C^\times$, its adjacency eigenvalue is
\[
 \lambda_\chi=\sum_O B_Ow_O,\qquad
 w_O=\sum_{s\in O}\chi(s)\in\mathbb R,
\]
where $B_O$ is a fair Bernoulli variable. Character orthogonality gives
$\mathbb E\lambda_\chi=-1/2$ for every nontrivial character. For the
trivial character the eigenvalue is the degree $d$, with expectation
$(N-1)/2$. The characters form an orthogonal eigenbasis, so all
nonprincipal eigenvalues are covered by this calculation.

Since $|w_O|\le|O|\le2$, we have $\sum_Ow_O^2\le2(N-1)\le2N$.
For the centered sum $Z=\sum_O(B_O-1/2)w_O$, independence and
$\cosh u\le\exp(u^2/2)$ yield
\[
 \mathbb E e^{tZ}=\prod_O\cosh(tw_O/2)
 \le\exp(t^2N/4).
\]
Markov's inequality followed by optimization in $t>0$ gives
$\Pr(|Z|\ge u)\le2\exp(-u^2/N)$. The same bound applies to the
trivial character. Taking $u=\sqrt{N\log(8N)}$ and a union bound over
all $N$ characters leaves failure probability at most $1/4$.
Consequently some inverse-closed $S$ satisfies simultaneously
\[
 |d-(N-1)/2|\le u,\qquad
 |\lambda_\chi+1/2|\le u\quad(\chi\ne1).
\]
For $N\ge1024$, these inequalities give $d\ge N/4$ and also
$N-1-d\ge N/4$. For example, $\log(8N)\le N/32$ throughout this
range, which immediately implies the two degree estimates.

\paragraph{From eigenvalues to both homogeneous-set bounds.}
Set $\Lambda=u+1/2$. For an independent set $I$ of size $a$, center
its indicator by $z=1_I-(a/N)1$. Since the graph is $d$-regular,
\[
 z^TAz=-da^2/N,\qquad z^Tz=a-a^2/N.
\]
Every nonprincipal eigenvalue is at least $-\Lambda$, so
$da^2/N\le\Lambda(a-a^2/N)$. For $a>0$, rearrangement gives
\[
 a\le\frac{N\Lambda}{d+\Lambda}\le4\Lambda.
\]
The complement is also a Cayley graph. Its nonprincipal eigenvalues
are $-1-\lambda_\chi$, which have the same bound $\Lambda$, and its
degree is at least $N/4$. Applying the independent-set estimate to
that graph bounds cliques in the original one. We have therefore proved
\[
 \max(\alpha,\omega)\le4\sqrt{N\log(8N)}+2
 \qquad(N\ge1024).
\]
Increasing an absolute constant handles smaller groups. This is a
uniform existence result, but it is much weaker than $C\log N$.

\paragraph{Why independent-edge counting fails.}
In $A=\mathbb F_2^m$, let $H$ be a $d$-dimensional subgroup. Its
$2^d$ vertices form a clique precisely when every nonzero element of
$H$ belongs to $S$. That event has probability
$2^{-(2^d-1)}$, rather than $2^{-\binom{2^d}{2}}$. The corresponding
independent-set event has the same probability. All edges with the
same difference are tied to a single random choice. This exact
calculation explains why applying the usual binomial random-graph
union bound to Cayley edges would be invalid. It does not alone
prove a high-probability obstruction: that would also require control
of overlaps among subgroups, as addressed by more substantial work
such as [R1].

\paragraph{Verification and final status.}
All estimates treat elements of order two as singleton orbits, and
complementation preserves the same identity-free convention. The
spectral bound uses an actual simultaneous event of positive probability,
not independently chosen graphs for the two homogeneous-set conditions.
\textbf{UNRESOLVED.} The first missing step is a construction or counting
argument reducing the square-root-scale bound to logarithmic size
uniformly over all abelian groups. The independent character estimates
proved here do not supply that improvement.

\subsection{Sources}
\begin{itemize}
\item[S1] ULAM.AI and the UnsolvedMath Contributors, supplied problems.json, record 3147 (OPG-372), OpenGarden collection. Catalogue identity and source transcription; CC BY 4.0.
\url{https://huggingface.co/datasets/ulamai/UnsolvedMath}.
\item[S2] Open Problem Garden, Ramsey properties of Cayley graphs. Original problem and discussion; the mathematical formulation below is expanded from the supplied source record.
\url{http://www.openproblemgarden.org/op/ramsey_properties_of_cayley_graphs}.

\item[R1] Ben Green, \emph{Counting sets with small sumset, and the
clique number of random Cayley graphs}, 2003. Primary abstract checked
1 October 2026; its sum-graph convention is distinguished in the text.
\url{https://arxiv.org/abs/math/0304183}.

\end{itemize}
\end{document}
